\documentclass[addpoints,spanish, 12pt,a4paper]{exam}
%\documentclass[answers, spanish, 12pt,a4paper]{exam}
\printanswers
\renewcommand*\half{.5}
\pointpoints{punto}{puntos}
\hpword{Puntos:}
\vpword{Puntos:}
\htword{Total}
\vtword{Total}
\hsword{Resultado:}
\hqword{Ejercicio:}
\vqword{Ejercicio:}

\usepackage[utf8]{inputenc}
\usepackage[spanish]{babel}
\usepackage{eurosym}
%\usepackage[spanish,es-lcroman, es-tabla, es-noshorthands]{babel}


\usepackage[margin=1in]{geometry}
\usepackage{amsmath,amssymb}
\usepackage{multicol}
\usepackage{yhmath}

\pointsinrightmargin % Para poner las puntuaciones a la derecha. Se puede cambiar. Si se comenta, sale a la izquierda.
\extrawidth{-2.4cm} %Un poquito más de margen por si ponemos textos largos.
\marginpointname{ \emph{\points}}

\usepackage{graphicx}

\graphicspath{{../img/}} 

\newcommand{\class}{1º Bachillerato CIE}
\newcommand{\examdate}{\today}
\newcommand{\examnum}{Parcial 1ªEv.}
\newcommand{\tipo}{A}


\newcommand{\timelimit}{45 minutos}

\renewcommand{\solutiontitle}{\noindent\textbf{Solución:}\enspace}


\pagestyle{head}
\firstpageheader{\includegraphics[width=0.2\columnwidth]{header_left}}{\textbf{Departamento de Matemáticas\linebreak \class}\linebreak \examnum}{\includegraphics[width=0.1\columnwidth]{header_right}}
\runningheader{\class}{\examnum}{Página \thepage\ de \numpages}
\runningheadrule


\usepackage{pgf,tikz,pgfplots}
\pgfplotsset{compat=1.15}
\usepackage{mathrsfs}
\usetikzlibrary{arrows}


\begin{document}

\noindent
\begin{tabular*}{\textwidth}{l @{\extracolsep{\fill}} r @{\extracolsep{6pt}} }
\textbf{Nombre:} \makebox[3.5in]{\hrulefill} & \textbf{Fecha:}\makebox[1in]{\hrulefill} \\
 & \\
\textbf{Tiempo: \timelimit} & Tipo: \tipo 
\end{tabular*}
\rule[2ex]{\textwidth}{2pt}
Esta prueba tiene \numquestions\ ejercicios. La puntuación máxima es de \numpoints. 
La nota final de la prueba será la parte proporcional de la puntuación obtenida sobre la puntuación máxima. 

\begin{center}


\addpoints
 %\gradetable[h][questions]
	\pointtable[h][questions]
\end{center}

\noindent
\rule[2ex]{\textwidth}{2pt}

\begin{questions}

%\question 
%
%\begin{parts}
%\part[2] 
%\begin{solution}
%\end{solution}
%
%
%\end{parts}
%\addpoints

% \question[1] Indica a cuáles de los conjuntos
% $\mathbb{N}$, $\mathbb{Z}$, $\mathbb{Q}$, $\mathbb{R}$ pertenecen cada uno de los siguientes números:
% \begin{center}
% \begin{tabular}{|c |c |c |c |c|}\hline
% &$\mathbb{N}$& $\mathbb{Z}$& $\mathbb{Q}$&$\mathbb{R}$\\ 
% \hline
% $\frac{8}{16}$&&&&\\
% \hline
% $\sqrt[3]{-27}$&&&&\\
% \hline
% $3.0\wideparen{1}$&&&&\\
% \hline
% $-\frac{12}{4}$&&&&\\
% \hline
% $-\sqrt{25}$&&&&\\
% \hline
% $\sqrt{8}$&&&&\\
% \hline
% $4$&&&&\\
% \hline
% $\pi$&&&&\\
% \hline
% $\sqrt{-4}$&&&&\\
% \hline
% $\frac{39}{13}$&&&&\\
% \hline
% \end{tabular}

% \end{center}

%         \question[2]  Dados los conjuntos:\\
%          $ A=\left\{ x \in \mathbb{R}| 6 \leq x  < 8 \right\}, \\ B=\left(-\infty, -3\right) \cup \left(3, \infty\right)  \\  C=\left\{ x \in \mathbb{R}| \left|{x - 3}\right|\leq12 \right\}  $
         
%          \begin{parts}
%           % \part Indica razonadamente, si existen, los supremos, ínfimos, máximos y mínimos de los conjuntos $A$, $B$ y $C$
%           \part Calcula $A \cup  B$ , $A \cap B$ y $(A \cup B) \cap C$ , 
% y expresa los resultados en forma de Intervalos.
%          \end{parts}
         
% \begin{solution}  $ C=\left[-9, 15\right] \ \ A \cup  B = \left(-\infty, -3\right) \cup \left(3, \infty\right)  \\  A \cap B= \left[6, 8\right)   \\  (A \cup B) \cap C= \left[-9, -3\right) \cup \left(3, 15\right] $  \end{solution}

       
        
% \question[1\half] Halla el término general de las siguientes sucesiones. Indica cuáles de ellas son progresiones aritméticas y cuáles progresiones geométricas:
% \begin{parts}
%     \part \( 4, 9, 14, 19, 24, 29, \ldots \)
%     \part \( 2, -6, 18, -54, 162, \ldots \)
% \end{parts}
% \begin{solution}
% $\text{PA: 1. } a_n = 5n - 1; \, \text{SG: -3. } a_n = (-3)^{n-1}\cdot 2$\end{solution}


\question Resuelve las siguientes ecuaciones:

\begin{parts}
%     \part[2]$2 \sqrt{x + 4} - \sqrt{x - 1} = 4$
%     \begin{solution}
% 1. Aislamos uno de los radicales:
%    \[
%    2 \sqrt{x + 4} = \sqrt{x - 1} + 4
%    \]
% 2. Elevamos ambos lados al cuadrado:
%    \[
%    4(x + 4) = (\sqrt{x - 1} + 4)^2
%    \]
%    \[
%    4x + 16 = x - 1 + 8 \sqrt{x - 1} + 16
%    \]
% 3. Simplificamos y resolvemos la ecuación resultante.
% $3x+1=8 \sqrt{x - 1} \to x=13/9, x=5
%  $

% \end{solution}
%     \part[1\half] $\sqrt{2x+3}-\sqrt{x+2}=2$
%     \begin{solution}
%     $2x+3=4+4\sqrt{x+2}+x+2 \to x-3=4\sqrt{x+2} \to x^2-6x+9=16x+32\to x^2-22x-23=0$\\
%     $x=23$ sí, $x=-1$ No
% \end{solution}
%     \part[1\half] $\sqrt[3]{x + 6} = x$ \begin{solution}
% 1. Elevamos ambos lados al cubo:
%    \[
%    x + 6 = x^3
%    \]
% 2. Reorganizamos la ecuación:
%    \[
%    x^3 - x - 6 = 0
%    \]
% 3. Intentamos con divisores de 6 para encontrar posibles raíces.
% $$x^3 - x - 6=\left(x - 2\right) \left(x^{2} + 2 x + 3\right)\to x=2$$
% \end{solution}
%     \part[1\half] $9^x - 6 \cdot {3^{x + 1}} + 81 = 0 $
% \begin{solution}
%     $\rightarrow \left [ 2\right ]$
% \end{solution}
\part[1]  $ 2 \log x - \log (x + 6) = 3 \log 2 $  \begin{solution}  $ \left [ 12\right ] $  \end{solution}
    \part[1] $|2x-4|=|x+2|$ \begin{solution} $\left[ \frac{2}{3}, \  6\right]$\end{solution}
\end{parts}


% \question[1] Da una cota del error absoluto y otra del error relativo cometidos al aproximar el número aúreo $\left(\phi\right)$ a $1.62$

        \question[3] Resuelve, justificadamente, los siguientes sistemas de inecuaciones:
        \begin{multicols}{2}
        \begin{parts} 
         \part  $$ \left\{\begin{matrix}{( {x - 1} )^2} - {( {x + 3} )^2} \leq 0\\x - 3( {x - 1} ) \geq 3 \end{matrix}\right. $$  
         \begin{solution}  $ \left[-1, 0\right] $  \end{solution}
         % \part[2]  $$ \left\{\begin{matrix}x \geq 0 \\0 \leq y \leq 3\\x - 2y \leq 10\\x + y \geq 10\\\end{matrix}\right. $$  \begin{solution}    \end{solution}  

         \part  $$ \left\{\begin{matrix}x \geq 0 \\0 \leq y \leq 3 \\x + y \leq 5\\\end{matrix}\right. $$  \begin{solution}    \end{solution} 
        \end{parts}
                  

        
        % \question Calcula expresando el resultado en forma de fracción algebraica irreducible:
        % \begin{multicols}{1}
        % \begin{parts} \part[1]  $ \dfrac{2+{ \dfrac{1}{x}}}{{2 + \dfrac{1}{{1 + \dfrac{1}{x}}}}} $  \begin{solution}  $ \dfrac{2 x^{2} + 3 x + 1}{3 x^{2} + 2 x} $  \end{solution}
        % \end{parts}
        % \end{multicols}
            
        \end{multicols}

% \question[1] Resolver la ecuación
% \[
% 4x^3 - 6x^2 - x = -\frac{3}{2},
% \]
% sabiendo que una solución es \( \frac{1}{2} \).

% \begin{solution}
% 1. Llevamos todos los términos a un lado:
%    \[
%    4x^3 - 6x^2 - x + \frac{3}{2} = 0
%    \]
% 2. Multiplicamos por 2 para eliminar fracciones:
%    \[
%    8x^3 - 12x^2 - 2x + 3 = 0
%    \]
% 3. Factorizando el polinomio 
% $\frac{\left(2 x - 3\right) \left(2 x - 1\right) \left(2 x + 1\right)}{2} \to \left[ - \frac{1}{2}, \  \frac{1}{2}, \  \frac{3}{2}\right]
% $
% \end{solution}

% \question[1] Resolver la ecuación
% \[
% 2x^3 - 3x^2 = -\frac{1}{2},
% \]
% sabiendo que una de sus raíces es \( \frac{1}{2} \).

% \begin{solution}
% 1. Llevamos todos los términos a un lado:
%    \[
%    2x^3 - 3x^2 + \frac{1}{2} = 0
%    \]
% 2. Multiplicamos por 2 para eliminar fracciones:
%    \[
%    4x^3 - 6x^2 + 1 = 0
%    \]
% 3. $$\left(2 x - 1\right) \left(2 x^{2} - 2 x - 1\right) \to \left[ \frac{1}{2}, \  \frac{1}{2} - \frac{\sqrt{3}}{2}, \  \frac{1}{2} + \frac{\sqrt{3}}{2}\right]$$
% \end{solution}

% \question Resuelve la siguiente ecuación:
% $9^x - 6 \cdot {3^{x + 1}} + 81 = 0 $
% \begin{solution}
%     $\rightarrow \left [ 2\right ]$
% \end{solution}


% \question Resuelve los siguientes sistemas:
% \begin{parts}
%     \part[1\half] $\left\{\begin{matrix}{2^x} + {2^y} = 24\\ \\{2^{x + y}} = 128\end{matrix}\right. $

% \begin{solution}
%     $\rightarrow \left [ \left \{ x : 3, \quad y : 4\right \}, \quad \left \{ x : 4, \quad y : 3\right \}\right ]$
% \end{solution}

%     \part[1\half] $\left\{\begin{matrix}\log_{2}(x) + \log_{2}(x + y) = 4\\ \\x + y = 2 \end{matrix}\right. $
% \begin{solution}
%     $\rightarrow \left [ \left \{ x : 8, \quad y : -6\right \}\right ]$
% \end{solution}

% \end{parts}


% \question Resuelve:
% $$\left\{\begin{matrix}3\log x - 2\log y = 10\\\log x + 3\log y = 7\end{matrix}\right. $$

% \begin{solution}
%     $\rightarrow \left [ \left \{ x : 10000, \quad y : 10\right \}\right ]$
% \end{solution}

% \question[2] Resuelve el siguiente sistema de ecuaciones:
% \[
% \begin{cases}
% \dfrac{1}{x} + y = 3 \\ \\
% \dfrac{1}{x} - \dfrac{1}{y} = \dfrac{1}{2}
% \end{cases}
% \]

\begin{solution}
1. Llamamos \(u = \frac{1}{x}\) y \(v = y\), para obtener el sistema:
   \[
   u + v = 3
   \]
   \[
   u - \frac{1}{v} = \frac{1}{2}
   \]
2. Resolvemos este sistema para \(u\) y \(v\), y luego encontramos \(x\) y \(y\) a partir de \(u\) y \(v\).
\end{solution}

% \question Resuelve las siguientes ecuaciones:
% \begin{parts}
%     \part $|2x-4|=|x+2|$ \begin{solution} $\left[ \frac{2}{3}, \  6\right]$\end{solution}
% \end{parts}


% https://lareinadelasciencias.wordpress.com/wp-content/uploads/2014/01/problemas-ecuaciones-sistemas-no-lineales.pdf

%https://matematicasiesoja.wordpress.com/wp-content/uploads/2013/10/problemas-_metodo_gauss2.pdf

\question Resuelve mediante expresiones algebraicas y, en caso que se pueda, por Gauss:
\begin{parts}
    \part[1] Un padre ha comprado un jersey para cada uno de sus cinco hijos, gastándose en total
108,75 euros. Tres de los jerseys tenían un 15\% de descuento, y otro de ellos tenía un 20\%
de descuento. Sabiendo que inicialmente costaban lo mismo, ¿cuánto ha tenido que pagar
por cada jersey? 
\begin{solution}$3\cdot0.85x+0.8x+x=108.75 \to x=25 \to 21.25; 20; 25$\end{solution}

    \part[1] El área de un triángulo es 78 $cm^2$  y entre la base y la altura suman 25 $cm$. Calcula la base y la altura.
% \begin{solution}  $ \left\{\begin{matrix}\frac{xy}{2}=78\\ x+y=25\\ \end{matrix}\right.  \rightarrow  \\\left[\begin{matrix}1 & 1 & 25\\0 & \frac{x}{2} - \frac{y}{2} & - \frac{25 y}{2} + 78\end{matrix}\right] \rightarrow  \left [ \left \{ x : 12, \quad y : 13\right \}, \quad \left \{ x : 13, \quad y : 12\right \}\right ] $  \end{solution}

%     \part[1] Miguel y Ana tienen un perro. Averigua el peso de cada uno de los tres sabiendo que Miguel y Ana pesan 50
% kg juntos, y Ana y su perro 29 kg y, finalmente, Miguel y el perro 35 kg.  \begin{solution}  $ \left\{\begin{matrix}x+y = 50\\ y+z=29\\ x+z=35\\ \end{matrix}\right.  \rightarrow  \\\left[\begin{matrix}1 & 1 & 0 & 50\\0 & 1 & 1 & 29\\0 & 0 & 2 & 14\end{matrix}\right] \rightarrow  \left \{ x : 28, \quad y : 22, \quad z : 7\right \} $  \end{solution}
  
%     \part[2] En una residencia de estudiantes se compran semanalmente 110 helados de distintos
% sabores: vainilla, chocolate y nata. El presupuesto destinado para esta compra es de 540
% euros y el precio de cada helado es de 4 euros el de vainilla, 5 euros el de chocolate y 6
% euros el de nata. Conocidos los gustos de los estudiante, se sabe que entre helados de
% chocolate y de nata se han de comprar el 20\% más que de vainilla 
% \begin{solution}    
% \textbf{Resolución del problema:}

% Sea:
% \[
% x: \text{número de helados de vainilla}, \quad
% y: \text{número de helados de chocolate}, \quad
% z: \text{número de helados de nata}.
% \]

% Planteamos las ecuaciones:
% \[
% x + y + z = 110, \quad 4x + 5y + 6z = 540, \quad y + z = 1.2x.
% \]

% \textbf{Sistema matricial:}
% \[
% \begin{bmatrix}
% 1 & 1 & 1 \\ 
% 4 & 5 & 6 \\ 
% -1.2 & 1 & 1
% \end{bmatrix}
% \begin{bmatrix}
% x \\ 
% y \\ 
% z
% \end{bmatrix}
% =
% \begin{bmatrix}
% 110 \\ 
% 540 \\ 
% 0
% \end{bmatrix}.
% \]

% \textbf{Eliminación de Gauss:}

% Partimos de la matriz ampliada:
% \[
% \left[
% \begin{array}{ccc|c}
% 1 & 1 & 1 & 110 \\ 
% 4 & 5 & 6 & 540 \\ 
% -1.2 & 1 & 1 & 0
% \end{array}
% \right].
% \]

% Tras operaciones elementales:
% \[
% \left[
% \begin{array}{ccc|c}
% 1 & 1 & 1 & 110 \\ 
% 0 & 1 & 2 & 100 \\ 
% 0 & 0 & 1 & 40
% \end{array}
% \right].
% \]

% \textbf{Sustitución hacia atrás:}
% \[
% z = 40, \quad y + 2z = 100 \implies y = 20, \quad x + y + z = 110 \implies x = 50.
% \]

% \textbf{Solución final:}
% \[
% x = 50 \quad \text{(vainilla)}, \quad y = 20 \quad \text{(chocolate)}, \quad z = 40 \quad \text{(nata)}.
% \]
% \end{solution}


\part[2] Una empresa envasadora ha comprado un total de 1500 cajas de pescado en tres mercados 
diferentes, a un precio por caja 30, 20 y 40  euros respectivamente. el coste total de la operación ha sido 
de 40500\euro. Calcular cuánto ha pagado la empresa en cada mercado, sabiendo que en el primero de ellos ha 
comprado el 30\% de las cajas

\begin{solution}
    Llamamos $x$, $y$ y $z$ cajas de 30, 20 y 40\euro \  respectivamente.
    El sistema a resolver es:$\left\{ \begin{matrix}x + y + z = 1500 \\ 30 x + 20 y + 40 z = 40500 \\ x = 450 \\ \end{matrix}\right.$ \\ 
    \textbf{Discusión y resolución por Gauss:} Escalonando la matriz ampliada tenemos\\$A^*= \left(\begin{matrix}1 & 1 & 1 & 1500\\30 & 20 & 40 & 40500\\1 & 0 & 0 & 450\end{matrix}\right) \thicksim \left(\begin{matrix}1 & 1 & 1 & 1500\\0 & -10 & 10 & -4500\\0 & 0 & -2 & -600\end{matrix}\right)$. \\  De los valores de la última fila podemos concluir:\begin{itemize}\item S.C.D.\begin{itemize}\item $\left(\begin{matrix}0 & 0 & -2 & -600\end{matrix}\right) \to z = 300$\end{itemize}\begin{itemize}\item $\left(\begin{matrix}0 & -10 & 10 & -4500\end{matrix}\right) \to y = 750$\end{itemize}\begin{itemize}\item $\left(\begin{matrix}1 & 1 & 1 & 1500\end{matrix}\right) \to x = 450$\end{itemize}\end{itemize}  \textbf{Por rangos y determinantes:} \\$\left|A\right|=\left|\begin{matrix}1 & 1 & 1\\30 & 20 & 40\\1 & 0 & 0\end{matrix}\right|=20  \neq 0 $\begin{itemize} \item $rg(A)=3 \land rg(A^*)=3 \to $ S.C.D.  \\ \\ Por Cramer: \begin{itemize}\item $x=\frac{\left|\begin{matrix}1500 & 1 & 1\\40500 & 20 & 40\\450 & 0 & 0\end{matrix}\right|}{20}=\frac{9000}{20}=450$\item $y=\frac{\left|\begin{matrix}1 & 1500 & 1\\30 & 40500 & 40\\1 & 450 & 0\end{matrix}\right|}{20}=\frac{15000}{20}=750$\item $z=\frac{\left|\begin{matrix}1 & 1 & 1500\\30 & 20 & 40500\\1 & 0 & 450\end{matrix}\right|}{20}=\frac{6000}{20}=300$\end{itemize}\end{itemize}
    \textbf{SOLUCIÓN: } 13500 (450x30), 15000 y 12000 \euro.
\end{solution}

    
%     \part[1] En una reunión hay 22 personas, entre hombres, mujeres y niños. El doble del
% número de mujeres más el triple del número de niños, es igual al doble del número de 
% hombres. 
% \begin{solution}
%     Sea:
% \[
% x: \text{número de hombres}, \quad
% y: \text{número de mujeres}, \quad
% z: \text{número de niños}.
% \]

% Planteamos las ecuaciones:
% \[
% x + y + z = 22, \quad 2y + 3z = 2x.
% \]

% \textbf{Sistema matricial:}
% \[
% \begin{bmatrix}
% 1 & 1 & 1 \\ 
% -2 & 2 & 3
% \end{bmatrix}
% \begin{bmatrix}
% x \\ 
% y \\ 
% z
% \end{bmatrix}
% =
% \begin{bmatrix}
% 22 \\ 
% 0
% \end{bmatrix}.
% \]

% \textbf{Eliminación de Gauss:}

% Partimos de la matriz ampliada:
% \[
% \left[
% \begin{array}{ccc|c}
% 1 & 1 & 1 & 22 \\ 
% -2 & 2 & 3 & 0
% \end{array}
% \right].
% \]

% Operaciones elementales eliminan el coeficiente de \(x\) en la segunda fila:
% \[
% \left[
% \begin{array}{ccc|c}
% 1 & 1 & 1 & 22 \\ 
% 0 & 4 & 5 & 44
% \end{array}
% \right].
% \]

% Simplificamos la segunda fila:
% \[
% \left[
% \begin{array}{ccc|c}
% 1 & 1 & 1 & 22 \\ 
% 0 & 1 & \frac{5}{4} & 11
% \end{array}
% \right].
% \]

% \textbf{Sustitución hacia atrás:}

% De la segunda fila:
% \[
% y + \frac{5}{4}z = 11 \quad \implies \quad y = 11 - \frac{5}{4}z.
% \]

% De la primera fila:
% \[
% x + y + z = 22 \quad \implies \quad x = 22 - y - z.
% \]

% Sustituyendo \(y\):
% \[
% x = 22 - \left(11 - \frac{5}{4}z\right) - z \quad \implies \quad x = 11 - \frac{1}{4}z.
% \]

% Para que \(x, y, z\) sean números enteros, \(z = 4\):
% \[
% y = 11 - \frac{5}{4}(4) = 6, \quad x = 11 - \frac{1}{4}(4) = 10.
% \]

% \textbf{Solución final:}
% \[
% x = 10 \quad \text{(hombres)}, \quad y = 6 \quad \text{(mujeres)}, \quad z = 4 \quad \text{(niños)}.
% \]    
% \end{solution}


\end{parts}


% \question Resuelve los siguientes sistemas de ecuaciones no lineales

% \begin{parts}

% \part[1] 
% \[
% \begin{cases}
% 2x - y = 1 \\
% 3x - 2y^2 = 1
% \end{cases}
% \]

% \begin{solution}
% $x = \frac{5}{4}, \, y = \frac{3}{2}$ \quad y \quad $x = -\frac{1}{4}, \, y = \frac{1}{2}$
% \end{solution}

% \part[1] 
% \[
% \begin{cases}
% x^2 - 2y^2 = 1 \\
% x \cdot y = 6
% \end{cases}
% \]

% \begin{solution}
% $x = 3, \, y = 2$ \quad y \quad $x = -3, \, y = -2$
% \end{solution}

% \part[1] 
% \[
% \begin{cases}
% \sqrt{x + y} + y = 1 \\
% 2x + y = 2
% \end{cases}
% \]

% \begin{solution}
% $x = 1, \, y = 0$
% \end{solution}

% \part[1] 
% \[
% \begin{cases}
% 2x^2 + 3y^2 = 21 \\
% -5x^2 + 2y^2 = -43
% \end{cases}
% \]

% \begin{solution}
% $x = \pm 2, \, y = \pm 3$
% \end{solution}

% \part[1] 
% \[
% \begin{cases}
% \frac{y}{x} - \frac{3}{y} = \frac{1}{2} \\
% x - 3y = -5
% \end{cases}
% \]

% \begin{solution}
% $x = 1, \, y = 2$
% \end{solution}

% \end{parts}

% \question[2] Resuelve por Gauss el siguiente sistema de ecuaciones:
% \[
% \begin{cases}
% 3x + 2y - 2z = 4 \\
% 4x + y - z = 7 \\
% x + 4y - 4z = 0
% \end{cases}
% \]

% \begin{solution}
% $\left[\begin{matrix}3 & 2 & -2 & 4\\0 & - \frac{5}{3} & \frac{5}{3} & \frac{5}{3}\\0 & 0 & 0 & 2\end{matrix}\right]$ Luego S.I.
% \end{solution}



\question Resuelve las siguientes inecuaciones:
\begin{parts}
            % \part[1]  $ \dfrac{{{x^3} - 5{x^2} + 2x + 8}}{1-x^2} \leq 0 $  \begin{solution}  $  \left(1, 2\right] \cup \left[4, \infty\right) $  \end{solution}
        \part[1] $1-\dfrac{{x + 3}}{{x + 6}} \geq 0 $ \begin{solution} $\rightarrow \left(-6, \infty\right)$\end{solution}
\end{parts}



        

\addpoints

\end{questions}

\end{document}
\grid
